\documentclass[10pt,a4paper,fontset=fandol]{ctexart}
\usepackage[margin=2cm]{geometry}
\usepackage{needspace}
\ctexset{section={format=\large\sffamily\bfseries,beforeskip=14pt,afterskip=7pt},subsection={format=\normalsize\sffamily\bfseries,beforeskip=10pt,afterskip=5pt}}
\usepackage{fontspec,booktabs,longtable,array,xcolor,hyperref}
\setmainfont{TeX Gyre Pagella}
\setsansfont{TeX Gyre Heros}
\definecolor{reportblue}{HTML}{244D65}
\hypersetup{colorlinks=true,urlcolor=reportblue,linkcolor=reportblue}
\setlength{\parindent}{0pt}
\setlength{\parskip}{5pt}
\renewcommand{\arraystretch}{1.18}
\setlength{\tabcolsep}{5pt}
\emergencystretch=2em
\urlstyle{same}
\pagestyle{plain}
\title{\sffamily Jev 在 MoralChoice 中的道德选择与稳定性}
\author{}
\date{2026 年 9 月 30 日}
\begin{document}
\begin{center}{\sffamily\LARGE Jev 在 MoralChoice 中的道德选择与稳定性\par}\vspace{8pt}{\small 2026 年 9 月 30 日}\end{center}\vspace{6pt}
\section*{摘要}
在 MoralChoice 的 1,367 个英语情境中，每个情境包含六个结构化判断。687 个低歧义情境的六种变体全部选择参考行动。680 个高歧义情境中，69.07\% 的判断选择 action1，90.59\% 的情境在六种变体下选择同一行动。高歧义情境的选择通常稳定，但对提问形式和顺序仍有敏感性。
\section{数据集}
发布数据包含 687 个低歧义和 680 个高歧义情境。低歧义情境以 action1 为参考行动，高歧义情境不设统一正确答案。每个情境包括 ab、repeat、compare 三种形式及各自的正逆序，共 8,202 个判断。模型版本为 typesafe/jev-1.13-20260917。

形式和顺序沿用原研究设计，输出改为 Jev 的结构化二选一。同一情境六个问题在一次请求中回答；repeat 不再测试逐字文本复述。顺序及跨形式稳定性为本次适配的分析指标，六个判断不作为六个独立情境。

\section{结果}
\Needspace{8\baselineskip}
\subsection{整体结果}
低歧义以发布数据的参考行动为判据，随机二选一参照为 50\%。高歧义仅描述选择分布，不计算正确率。

{\small
\begin{longtable}{@{}>{\raggedright\arraybackslash}p{\dimexpr 0.2400\textwidth-2\tabcolsep\relax}>{\raggedright\arraybackslash}p{\dimexpr 0.3800\textwidth-2\tabcolsep\relax}>{\raggedright\arraybackslash}p{\dimexpr 0.3800\textwidth-2\tabcolsep\relax}@{}}
\toprule
\textbf{指标} & \textbf{低歧义} & \textbf{高歧义}\\
\midrule
\endfirsthead
\toprule
\textbf{指标} & \textbf{低歧义} & \textbf{高歧义}\\
\midrule
\endhead
情境数 & 687 & 680\\
判断数 & 4122 & 4080\\
action1 选择比例 & 100.00\% & 69.07\%\\
六种变体选择同一行动 & 100.00\% & 90.59\%\\
平均 P(action1) & 99.82\% & 67.11\%\\
\bottomrule
\end{longtable}
}
低歧义的情境级全部选择参考行动比例为 100\%，Wilson 95\% 区间为 99.44\%–100\%。高歧义的 action1 比例受发布数据中行动内容的排列影响，不等于首位选项偏好。

\Needspace{8\baselineskip}
\subsection{高歧义稳定性}
{\small
\begin{longtable}{@{}>{\raggedright\arraybackslash}p{\dimexpr 0.2400\textwidth-2\tabcolsep\relax}>{\raggedright\arraybackslash}p{\dimexpr 0.1900\textwidth-2\tabcolsep\relax}>{\raggedright\arraybackslash}p{\dimexpr 0.1900\textwidth-2\tabcolsep\relax}>{\raggedright\arraybackslash}p{\dimexpr 0.1900\textwidth-2\tabcolsep\relax}>{\raggedright\arraybackslash}p{\dimexpr 0.1900\textwidth-2\tabcolsep\relax}@{}}
\toprule
\textbf{形式} & \textbf{正逆序一致} & \textbf{一致率} & \textbf{平均概率变化} & \textbf{首位选择率}\\
\midrule
\endfirsthead
\toprule
\textbf{形式} & \textbf{正逆序一致} & \textbf{一致率} & \textbf{平均概率变化} & \textbf{首位选择率}\\
\midrule
\endhead
ab & 661 / 680 & 97.21\% & 3.15\% & 49.04\%\\
repeat & 658 / 680 & 96.76\% & 3.25\% & 48.97\%\\
compare & 643 / 680 & 94.56\% & 6.36\% & 47.57\%\\
\bottomrule
\end{longtable}
}
616/680 个情境的六种变体选择相同。compare 的正逆序一致率最低，表明是／否比较形式更容易随顺序改变底层选择。六种概率的情境内平均极差为 10.75 个百分点。

\Needspace{8\baselineskip}
\subsection{辅助规则分组}
辅助规则标签来自发布数据，并可能相互重叠。仅比较一个行动被标为违反、另一个明确标为不违反该规则的情境。完整聚合表保存在机器可读汇总中。

{\small
\begin{longtable}{@{}>{\raggedright\arraybackslash}p{\dimexpr 0.3200\textwidth-2\tabcolsep\relax}>{\raggedright\arraybackslash}p{\dimexpr 0.2267\textwidth-2\tabcolsep\relax}>{\raggedright\arraybackslash}p{\dimexpr 0.2267\textwidth-2\tabcolsep\relax}>{\raggedright\arraybackslash}p{\dimexpr 0.2267\textwidth-2\tabcolsep\relax}@{}}
\toprule
\textbf{规则标签} & \textbf{高歧义情境数} & \textbf{选择不违反该规则} & \textbf{六种变体均不违反}\\
\midrule
\endfirsthead
\toprule
\textbf{规则标签} & \textbf{高歧义情境数} & \textbf{选择不违反该规则} & \textbf{六种变体均不违反}\\
\midrule
\endhead
break\_law & 195 & 84.02\% & 76.92\%\\
break\_promise & 108 & 68.21\% & 65.74\%\\
cheat & 138 & 93.24\% & 91.30\%\\
death & 84 & 64.88\% & 59.52\%\\
deceive & 262 & 86.39\% & 82.44\%\\
disable & 144 & 67.13\% & 61.11\%\\
duty & 409 & 86.15\% & 81.66\%\\
freedom & 219 & 54.26\% & 48.86\%\\
pain & 356 & 57.12\% & 51.97\%\\
pleasure & 276 & 50.60\% & 45.65\%\\
\bottomrule
\end{longtable}
}
\Needspace{8\baselineskip}
\subsection{调用成本}
{\small
\begin{longtable}{@{}>{\raggedright\arraybackslash}p{\dimexpr 0.2400\textwidth-2\tabcolsep\relax}>{\raggedright\arraybackslash}p{\dimexpr 0.3800\textwidth-2\tabcolsep\relax}>{\raggedright\arraybackslash}p{\dimexpr 0.3800\textwidth-2\tabcolsep\relax}@{}}
\toprule
\textbf{输入 token} & \textbf{输出 token} & \textbf{已报告费用（美元）}\\
\midrule
\endfirsthead
\toprule
\textbf{输入 token} & \textbf{输出 token} & \textbf{已报告费用（美元）}\\
\midrule
\endhead
1,098,374 & 250,161 & 0.046131708\\
\bottomrule
\end{longtable}
}
1,367 次请求全部成功，无重试或未知费用。

\Needspace{8\baselineskip}
\section{结论}
Jev 在低歧义情境下与参考选择完全一致。高歧义情境没有统一正确答案，约九成情境的底层行动在六种形式和顺序组合中保持一致；其余变化主要集中在比较形式。此结果描述结构化选择及其稳定性，不将集中的选择概率解释为道德正确性或人类意见比例。
\Needspace{5\baselineskip}
\section*{参考文献}
\begin{enumerate}
\item Scherrer et al. (2023). Evaluating the Moral Beliefs Encoded in LLMs. NeurIPS. \url{https://proceedings.neurips.cc/paper_files/paper/2023/file/a2cf225ba392627529efef14dc857e22-Paper-Conference.pdf}
\item MoralChoice data release, commit 89c0fe7b158b5ade5d10e0644c1aa20ab4c78cbe; CC BY 4.0. \url{https://huggingface.co/datasets/ninoscherrer/moralchoice}
\end{enumerate}
\end{document}
